% GENERATED FILE. Edit canonical lessons, then run npm run generate:books.
% canonical-source-hash: fnv1a64:e53b6dd9ed4e30ce
% canonical-lessons: ES-C429-instituto, ES-C429-nota, ES-C429-calificacion, ES-C429-aprobar, ES-C429-suspender, ES-C429-repaso-notas, ES-C429-sintesis-boletin

\chapter{The Instituto, And What A Mark Is Called}
\label{ch:instituto-y-nota}
\hlchaptermodality{429}

\begin{chapteropening}
\textbf{By the end of this chapter:} \emph{I can name the school a Spanish teenager goes to, read a report card, say what mark I got, and say whether I passed or failed.}

Read a three-row report card from a Spanish secondary school, say which subjects passed and which hangs, and answer the spoken version of the same question.
\end{chapteropening}

\section[el instituto]{el instituto — the building between the two you already know}

\label{lesson:ES-C429-instituto}



\begin{quote}
You can say \emph{la escuela} and \emph{la universidad}. A Spanish fifteen-year-old goes to neither. Where do they go?
\end{quote}

\subsection*{You'll want to know: el instituto}
\textbf{el instituto} — the secondary school, for students of about twelve to eighteen.

On a sign it is often shortened to \textbf{IES} — \emph{Instituto de Educación Secundaria}.

\begin{cousinweb}
Latin \textbf{instituere} — \emph{in-} plus \textbf{statuere}, \textquotedblleft{}to set up, to stand a thing upright.\textquotedblright{} An \textbf{institutum} was \textbf{something that had been set up}: a custom, an arrangement, a foundation. A Spanish \emph{instituto} is a \textbf{thing that was established}, and the school sense is one narrow use of a very wide word.

\emph{Statuere} is behind a crowd of words you can already half-read: \emph{estatua}, \emph{estado}, \emph{estatura}, \emph{constituir}. Every one of them is about \textbf{standing}.
\end{cousinweb}

\begin{grammarlens}[title={Grammar Lens: the false friend, and the ladder it breaks}]
English \emph{institute} names a research body — the Institute of Physics. Spanish \emph{instituto} can mean that too, but its everyday sense is \textbf{the school}, and that is the one a reading paper uses.

The three buildings in order:

\noindent
\begin{tabularx}{\linewidth}{@{}>{\raggedright\arraybackslash}X>{\raggedright\arraybackslash}X@{}}
\toprule
\textbf{age} & \textbf{building} \\
\midrule
about 6 to 12 & la escuela \\
about 12 to 18 & el instituto \\
after 18 & la universidad \\
\bottomrule
\end{tabularx}

Without the middle rung the ladder has a gap in exactly the years an exam candidate is most likely to be asked about.

The stress is \emph{ins-ti-TU-to}, on the second-last syllable, and no accent is written because a word ending in a vowel is stressed there by default.
\end{grammarlens}

\subsection*{Your turn}
Say these aloud:

\begin{itemize}
\raggedright
  \item \textquotedblleft{}el instituto\textquotedblright{} — \emph{ins-ti-TU-to}
  \item the three buildings in order of age
  \item what the letters \emph{IES} stand for
  \item what \emph{statuere} meant, and one Spanish word that still shows it
\end{itemize}

\begin{tcolorbox}[breakable,colback=teal!4,colframe=teal!35!black,title={Before you move on}]
Where does a fifteen-year-old go? (\emph{Al instituto}.) What does English \emph{institute} name instead? (A research body.) What was an \emph{institutum} in Latin? (Something that had been set up.)
\end{tcolorbox}
\section[la nota]{la nota — a mark, in the oldest sense of the word}

\label{lesson:ES-C429-nota}



\begin{quote}
You can say \emph{el examen}. Now try to say what you got for it, and notice you have the event and not the result.
\end{quote}

\subsection*{You'll want to know: la nota}
\textbf{la nota} — the mark, the grade.

\emph{¿Qué nota tienes?} — \textquotedblleft{}What mark did you get?\textquotedblright{} In Spain the answer is a number \textbf{out of ten}, and five is where passing begins.

\begin{cousinweb}
Latin \textbf{nota} meant \textbf{a mark scratched onto something} — a sign cut into wax, wood or skin so it could be told apart later.

Every modern sense grew out of that one scratch:

\noindent
\begin{tabularx}{\linewidth}{@{}>{\raggedright\arraybackslash}X>{\raggedright\arraybackslash}X@{}}
\toprule
\textbf{word} & \textbf{the mark in it} \\
\midrule
\textbf{nota} & the mark a teacher puts on your work \\
\textbf{notar} & to notice — to catch the mark \\
\textbf{notable} & worth marking \\
\bottomrule
\end{tabularx}

English \emph{note}, \emph{notice}, \emph{notable} and \emph{notorious} are the same word arriving by a different road.

So the grade sense is not a metaphor. \textbf{It is the literal one}, and a written note — a mark made on paper to be read later — is the extension.
\end{cousinweb}

\begin{grammarlens}[title={Grammar Lens: one word, two senses, and how to tell them apart}]
You have met \emph{nota} before, on a message left for somebody. Both senses are live, and what sorts them is the company the word keeps:

\noindent
\begin{tabularx}{\linewidth}{@{}>{\raggedright\arraybackslash}X>{\raggedright\arraybackslash}X@{}}
\toprule
\textbf{what stands nearby} & \textbf{which sense} \\
\midrule
\emph{examen}, \emph{clase}, a number & the mark \\
\emph{dejar}, \emph{escribir}, a person & the message \\
\bottomrule
\end{tabularx}

\emph{Sacar buena nota} is the set phrase for doing well — literally \textquotedblleft{}to take out a good mark,\textquotedblright{} the way you take out anything you have earned.
\end{grammarlens}

\subsection*{Your turn}
Say these aloud:

\begin{itemize}
\raggedright
  \item \textquotedblleft{}la nota\textquotedblright{} — the mark
  \item \textquotedblleft{}¿Qué nota tienes?\textquotedblright{}
  \item what a Latin \emph{nota} physically was
  \item the phrase for doing well in an exam (\emph{sacar buena nota})
\end{itemize}

\begin{tcolorbox}[breakable,colback=teal!4,colframe=teal!35!black,title={Before you move on}]
Out of how many is a Spanish school mark? (Ten.) Which sense of \emph{nota} came first, the mark or the message? (The mark — Latin \emph{nota} was a scratch.) What tells the two senses apart in a sentence? (The words standing next to it.)
\end{tcolorbox}
\section[la calificación]{la calificación — the word the paper uses when the person would say \emph{nota}}

\label{lesson:ES-C429-calificacion}



\begin{quote}
Say \emph{nota}. Now imagine it printed at the top of an official column on a school document, and expect a longer word.
\end{quote}

\subsection*{You'll want to know: la calificación}
\textbf{la calificación} — the grade, in its formal dress.

Spoken between students it is \emph{nota}. Printed on a report, a transcript or a form, it is \textbf{calificación}, and the two name the same thing.

\begin{cousinweb}
Medieval Latin \textbf{qualificare}, built from \textbf{qualis} (\textquotedblleft{}of what kind\textquotedblright{}) plus \textbf{facere} (\textquotedblleft{}to make\textquotedblright{}). To \emph{calificar} is \textbf{to say what kind a thing is}.

That is a startlingly honest description of what a grade does. A number out of ten is a claim about the \textbf{kind} of work in front of the reader, which is why the formal word for it is built from the question \emph{of what kind?}

English took the same compound as \textbf{qualify}, and kept the other half of the sense: to \emph{qualify} a statement is to say what kind of statement it is.
\end{cousinweb}

\begin{grammarlens}[title={Grammar Lens: where the accent sits, and why}]
\textbf{Calificación} ends in \textbf{-n}, which by default pulls the stress to the second-last syllable — \emph{ca-li-fi-CA-cion}. The stress is on the last one, so an accent has to be written: \textbf{-ción}.

Every Spanish word ending in \textbf{-ción} works this way, and you have met dozens:

\noindent
\begin{tabularx}{\linewidth}{@{}>{\raggedright\arraybackslash}X>{\raggedright\arraybackslash}X@{}}
\toprule
\textbf{word} & \textbf{stressed syllable} \\
\midrule
calificación & -ción \\
educación & -ción \\
información & -ción \\
\bottomrule
\end{tabularx}

In the plural the accent disappears — \emph{calificaciones} — because adding a syllable puts the stress where the default would have put it anyway.
\end{grammarlens}

\subsection*{Your turn}
Say these aloud:

\begin{itemize}
\raggedright
  \item \textquotedblleft{}la calificación\textquotedblright{} — \emph{ca-li-fi-ca-CIÓN}
  \item which of \emph{nota} and \emph{calificación} you would hear in a corridor
  \item what \emph{qualis} meant, and what a grade therefore claims
  \item the plural, and say why it loses its accent
\end{itemize}

\begin{tcolorbox}[breakable,colback=teal!4,colframe=teal!35!black,title={Before you move on}]
What question is hiding inside \emph{calificación}? (\emph{Of what kind?}) Where does its accent go, and where does it go in the plural? (On \emph{-ción}; the plural drops it.) Which English word is the same compound? (\emph{Qualify}.)
\end{tcolorbox}
\section[aprobar]{aprobar — to pass, and who is really doing the approving}

\label{lesson:ES-C429-aprobar}



\begin{quote}
You can name the exam and the mark. Say what you did with the exam when the mark was a five or better.
\end{quote}

\subsection*{You'll want to know: aprobar}
\textbf{aprobar} — to pass, as in to pass an exam.

\emph{He aprobado el examen.} — \textquotedblleft{}I have passed the exam.\textquotedblright{}

The mark that does it is a \textbf{five out of ten} and upward.

\begin{cousinweb}
Latin \textbf{approbare} — \emph{ad-} (\textquotedblleft{}towards\textquotedblright{}) plus \textbf{probare}, \textquotedblleft{}to test something and find it good.\textquotedblright{}

Look at the direction that gives you. English says \emph{I passed the exam}, with the student doing the acting and the exam standing still. The Latin behind \emph{aprobar} says something closer to \textbf{the examiner tested the work and found it good} — the judgement travels the other way.

\emph{Probare} is a generous root, and you can read its family off the page:

\noindent
\begin{tabularx}{\linewidth}{@{}>{\raggedright\arraybackslash}X>{\raggedright\arraybackslash}X@{}}
\toprule
\textbf{word} & \textbf{what it tests} \\
\midrule
\textbf{prob}ar & to try, to taste, to test \\
\textbf{prueb}a & a test, a piece of evidence \\
a\textbf{prob}ar & to test and find good \\
\bottomrule
\end{tabularx}

English \emph{probe}, \emph{proof}, \emph{prove} and \emph{approve} all come from the same verb.
\end{cousinweb}

\begin{grammarlens}[title={Grammar Lens: the vowel that breaks in the middle}]
\emph{Aprobar} is one of the verbs whose \textbf{o} splits into \textbf{ue} when the stress lands on it:

\noindent
\begin{tabularx}{\linewidth}{@{}>{\raggedright\arraybackslash}X>{\raggedright\arraybackslash}X@{}}
\toprule
\textbf{form} & \textbf{stem} \\
\midrule
yo apr\textbf{ue}bo & stressed \\
nosotros apr\textbf{o}bamos & unstressed \\
\bottomrule
\end{tabularx}

You have seen this before in \emph{poder} and \emph{contar}. The rule has not changed: the stem breaks where the stress falls and stays whole where it does not.
\end{grammarlens}

\subsection*{Your turn}
Say these aloud:

\begin{itemize}
\raggedright
  \item \textquotedblleft{}aprobar\textquotedblright{} — to pass
  \item \textquotedblleft{}He aprobado el examen\textquotedblright{}
  \item the \emph{yo} form, and hear where the stem breaks
  \item what \emph{probare} meant, and which direction the judgement travels
  \item the formal word a report prints for the number you got
\end{itemize}

\begin{tcolorbox}[breakable,colback=teal!4,colframe=teal!35!black,title={Before you move on}]
What mark passes? (Five out of ten.) What did \emph{probare} mean? (To test and find good.) What happens to the \textbf{o} in \emph{yo apruebo}? (It breaks into \textbf{ue}, because the stress lands on it.)
\end{tcolorbox}
\section[suspender]{suspender — to fail, and the picture the word is holding}

\label{lesson:ES-C429-suspender}



\begin{quote}
You can say \emph{aprobar}. Guess its opposite, and then read on to see what the Spanish word is actually picturing.
\end{quote}

\subsection*{You'll want to know: suspender}
\textbf{suspender} — to fail, as in to fail an exam or a subject.

\emph{He suspendido el examen.} — \textquotedblleft{}I have failed the exam.\textquotedblright{} A mark below \textbf{five} does it.

\begin{cousinweb}
Latin \textbf{suspendere} — \emph{sub-} (\textquotedblleft{}beneath\textquotedblright{}) plus \textbf{pendere}, \textquotedblleft{}to hang.\textquotedblright{} To suspend a thing is to leave it \textbf{hanging}.

That is not a dead metaphor here. A failed subject in Spain is \textbf{not finished}: it hangs over the student until a resit in a later term. There is even a word for one — \emph{una asignatura pendiente}, \textquotedblleft{}a hanging subject\textquotedblright{} — which people also use for anything in life they have been putting off.

\textbf{Pendere} built a shelf of words that all hang:

\noindent
\begin{tabularx}{\linewidth}{@{}>{\raggedright\arraybackslash}X>{\raggedright\arraybackslash}X@{}}
\toprule
\textbf{word} & \textbf{what hangs} \\
\midrule
\textbf{pend}iente & a slope, an earring, a thing outstanding \\
de\textbf{pend}er & to hang from something else \\
\textbf{pens}ar & to weigh a thing in the mind \\
\bottomrule
\end{tabularx}

English \emph{pendant}, \emph{pending}, \emph{suspense} and \emph{depend} hang from the same nail.
\end{cousinweb}

\begin{grammarlens}[title={Grammar Lens: the English false friend, and the Spanish one}]
English \emph{suspend} keeps only the harsher senses — a suspended player, a suspended sentence. The school sense does not survive the crossing, so an English speaker reading \emph{he suspendido} will hear something far more dramatic than what happened.

The pair, side by side:

\noindent
\begin{tabularx}{\linewidth}{@{}>{\raggedright\arraybackslash}X>{\raggedright\arraybackslash}X@{}}
\toprule
\textbf{mark} & \textbf{verb} \\
\midrule
five or above & aprobar \\
below five & suspender \\
\bottomrule
\end{tabularx}

Both verbs take the exam as their object, so the exam is what is passed or failed, rather than the student.
\end{grammarlens}

\subsection*{Your turn}
Say these aloud:

\begin{itemize}
\raggedright
  \item \textquotedblleft{}suspender\textquotedblright{} — to fail
  \item \textquotedblleft{}He suspendido el examen\textquotedblright{}
  \item the pair, with the mark that decides each
  \item what \emph{una asignatura pendiente} is, and what it means outside school
\end{itemize}

\begin{tcolorbox}[breakable,colback=teal!4,colframe=teal!35!black,title={Before you move on}]
What does \emph{pendere} mean? (To hang.) What is \emph{una asignatura pendiente}? (A subject left hanging, and by extension anything outstanding.) Which mark divides \emph{aprobar} from \emph{suspender}? (Five.)
\end{tcolorbox}
\section[el curso entero]{el curso entero — one school year, in five words}

\label{lesson:ES-C429-repaso-notas}



\begin{quote}
Name the building and the thing a teacher writes on your work, before you look at the table.
\end{quote}

\begin{grammarlens}[title={Grammar Lens: the year, in order}]
\noindent
\begin{tabularx}{\linewidth}{@{}>{\raggedright\arraybackslash}X>{\raggedright\arraybackslash}X@{}}
\toprule
\textbf{what happens} & \textbf{the word} \\
\midrule
you go to the building & el instituto \\
you sit in the room & la clase \\
you sit the thing & el examen \\
a number comes back & la nota \\
the document prints it & la calificación \\
five or above & aprobar \\
below five & suspender \\
\bottomrule
\end{tabularx}

Two of those seven were already yours. The chapter added the other five, and with them the track can follow a school year from the door to the report.
\end{grammarlens}

\begin{grammarlens}[title={Grammar Lens: two words for one number}]
\noindent
\begin{tabularx}{\linewidth}{@{}>{\raggedright\arraybackslash}X>{\raggedright\arraybackslash}X@{}}
\toprule
\textbf{where you meet it} & \textbf{which word} \\
\midrule
a friend asking in a corridor & la nota \\
a column heading on a document & la calificación \\
\bottomrule
\end{tabularx}

Same number, two registers. A reading paper prints the formal one; a listening paper says the short one.
\end{grammarlens}

\begin{grammarlens}[title={Grammar Lens: three Latin pictures}]
\noindent
\begin{tabularx}{\linewidth}{@{}>{\raggedright\arraybackslash}X>{\raggedright\arraybackslash}X@{}}
\toprule
\textbf{word} & \textbf{the picture under it} \\
\midrule
nota & a scratch cut into a surface \\
aprobar & a thing tested and found good \\
suspender & a thing left hanging \\
\bottomrule
\end{tabularx}

The third one is still literally true in Spain, where a failed subject hangs over a student until the resit.
\end{grammarlens}

\subsection*{Your turn}
Say these aloud:

\begin{itemize}
\raggedright
  \item the three school buildings by the age that goes to each
  \item the two words for the number, and where each one appears
  \item the two verbs, and the mark that divides them
  \item the picture hiding inside \emph{suspender}
\end{itemize}

\begin{tcolorbox}[breakable,colback=teal!4,colframe=teal!35!black,title={Before you move on}]
Which word heads a column on a report? (\emph{Calificación}.) Which verb is the Latin for \emph{found good}? (\emph{Aprobar}.) Which building takes a fifteen-year-old? (\emph{El instituto}.)
\end{tcolorbox}
\section[Practice]{Synthesis — reading a report card}

\label{lesson:ES-C429-sintesis-boletin}



\begin{quote}
Say the two words for a mark. One of them is about to appear as a column heading.
\end{quote}

\begin{tcolorbox}[breakable,skin=enhanced,colback=orange!6,colframe=orange!45!black,boxrule=0.5pt,arc=1mm,left=6pt,right=6pt,top=4pt,bottom=4pt,fonttitle=\bfseries,title={Read this}]
\begin{quote}
\textbf{IES Miguel Hernández} \textbf{Boletín de calificaciones}
\end{quote}

\noindent
\begin{tabularx}{\linewidth}{@{}>{\raggedright\arraybackslash}X>{\raggedright\arraybackslash}X@{}}
\toprule
\textbf{clase} & \textbf{calificación} \\
\midrule
Matemáticas & 7 \\
Historia & 4 \\
Inglés & 9 \\
\bottomrule
\end{tabularx}

Three rows, and everything in them is now readable. \textbf{IES} is the school. \textbf{Calificación} heads the column because this is a printed document. The numbers are out of ten.
\end{tcolorbox}

\subsection*{How to answer}
Two of the three rows passed and one did not. Say which verb goes with each:

\begin{quote}
\emph{He aprobado Matemáticas e Inglés.} \emph{He suspendido Historia.}
\end{quote}

The history row is the one that hangs. The student will carry it into the next term as \textbf{una asignatura pendiente}, and the family will ask about it in exactly those words.

\subsection*{The exchange}
\begin{quote}
— \emph{¿Qué notas has sacado?} — \emph{He aprobado casi todo. Historia no.}
\end{quote}

The question uses \textbf{notas} because it is spoken. The answer names the one that did not pass, which is what the question is really asking about.

\subsection*{Your turn}
Say these aloud:

\begin{itemize}
\raggedright
  \item what \emph{IES} at the top of the document stands for
  \item which subjects passed, using the right verb
  \item which subject hangs, and the phrase for it
  \item the spoken version of the question, and answer it about the 4
\end{itemize}

\begin{tcolorbox}[breakable,colback=teal!4,colframe=teal!35!black,title={Before you move on}]
Which word heads the column on the printed report? (\emph{Calificación}.) What happened to Historia? (\emph{Suspendida} — it hangs.) What is the spoken word for the same number? (\emph{La nota}.)
\end{tcolorbox}
